\begin{abstract}
	Este trabalho aborda o clássico Problema da Ponte e da Tocha, modelando-o
	como um Problema em Espaço de Estados e resolvendo-o por meio de quatro
	algoritmos de busca: Busca em Profundidade (DFS), Busca em Largura (BFS),
	Busca de Custo Mínimo (UCS) e Busca A*. São apresentadas a modelagem
	formal do problema, os detalhes de implementação de cada algoritmo e uma
	comparação experimental entre eles na instância clássica de quatro
	pessoas, considerando custo da solução, número de nós visitados e tempo
	de processamento. Os resultados mostram que apenas os algoritmos UCS e A*
	garantem a solução ótima.
\end{abstract}

\begin{IEEEkeywords}
	Fundamentos de IA, espaço de estado, busca em grafos, BFS, UFS, DFS, A*
\end{IEEEkeywords}
